%% 2-0 序章　「だんだん速く」を数で言うと？（第11回：X9 で題名と本文を変えた）
\makeatletter
\gdef\bk@partlabel{\R{第}{だい}2\R{巻}{かん}}\gdef\bk@partname{\R{数}{かず}と\R{量}{りょう}のことば}
%% ---- この章だけで使う絵の部品（頭文字 \bkiia…） ----
% 落ちるボール  \bkiiaball{(x,y)}{半径}{予想なら1}
\newcommand\bkiiaball[3]{\ifnum#3=1
    \draw[bkOrange, line width=0.9pt, dash pattern=on 2pt off 1.4pt, fill=bkOrange!12] #1 circle (#2);
  \else\shade[ball color=bkOrange] #1 circle (#2);\fi}
% 路線の1行  \bkiiaroute{y}{路線番号}{前の巻}{この巻の駅}{この先}
\newcommand\bkiiaroute[5]{%
  \draw[bkRoute#2, line width=3pt, line cap=round] (10,#1) -- (132,#1);
  \node[inner sep=0pt] at (4,#1) {\bkrouteicon[7mm]{#2}};
  \fill[white] (31,#1) circle (2.2);\draw[bkRoute#2, line width=1.1pt] (31,#1) circle (2.2);
  \node[anchor=north, font=\sffamily\footnotesize, text=bkGray, inner sep=0pt] at (31,#1-3.8) {#3};
  \fill[white, rounded corners=1.5mm] (50,#1-4.4) rectangle (100,#1+4.4);
  \draw[bkRoute#2, line width=1.6pt, rounded corners=1.5mm] (50,#1-4.4) rectangle (100,#1+4.4);
  \node[font=\sffamily\footnotesize, text=bkInk] at (75,#1) {#4};
  \fill[white] (119,#1) circle (2.2);\draw[bkRoute#2, line width=1.1pt] (119,#1) circle (2.2);
  \node[anchor=north, font=\sffamily\footnotesize, text=bkGray, inner sep=0pt] at (119,#1-3.8) {#5};}
% 読み方のカード  \bkiiacard{(x,y)}{アイコン}{色}{見出し}{中身}
\newcommand\bkiiacard[5]{%
  \node[anchor=north west, draw=#3, line width=1pt, fill=white, rounded corners=2.5mm,
    minimum width=66mm, minimum height=47mm, inner sep=0pt] at #1 {};
  \node[anchor=north west, inner sep=0pt] (ci) at ($#1+(3.5,-3)$) {\bkicon[1.5em]{#2}{#3}};
  \node[anchor=west, font=\sffamily\bfseries\small, text=#3!85!black, inner sep=0pt] at ([xshift=2mm]ci.east) {#4};
  \node[anchor=north west, text width=59mm, align=left, font=\small, inner sep=0pt] at ($#1+(3.5,-12)$) {#5};}
% 箇条書きの1項目（・で始め、2行目は字下げ）  \bkiiaitem{文}
\newcommand\bkiiaitem[1]{\par\noindent\hangindent=1\zw\hangafter=1 ・#1\par}
% 章番号の札  \bkiiapill{章}
\newcommand\bkiiapill[1]{\bk@pill{#1}}
% 会話の1行  \bkiiatalk{y}{sora|mina|delta}{名前}{せりふ}
\newcommand\bkiiatalk[4]{%
  \bkChar{#2}{(7,#1-6.5)}{12mm}
  \node[rectangle callout, callout absolute pointer={(15.5,#1)}, anchor=west, draw=bkTheme, line width=0.8pt,
    fill=white, rounded corners=2.5mm, inner xsep=3mm, inner ysep=1.8mm, text width=103mm, align=left,
    font=\small] at (21,#1) {#4};
  \node[anchor=north, font=\sffamily\scriptsize, text=bkGray, inner sep=0pt] at (7,#1-7.5) {#3};}
% 見出しのアイコン（路線：線と2つの駅）  \bksection{bkiiaroute}{…} で使う
\newcommand\bk@icon@bkiiaroute{%
  \draw[white, line width=0.13*\bk@iu, line cap=round] (-0.62,0) -- (0.62,0);
  \filldraw[fill=bkTheme, draw=white, line width=0.09*\bk@iu] (-0.36,0) circle (0.2) (0.36,0) circle (0.2);}
\makeatother

\bkchapter{2}{0}{「だんだん\R{速}{はや}く」を\\\R{数}{かず}で\R{言}{い}うと？}{\bklevel{1}}{「だんだん\R{速}{はや}く」を\R{数}{かず}で\R{言}{い}うと？}

\noindent \R{手}{て}に\R{持}{も}ったボールを、そっとはなします。ボールは、\R{下}{した}へ\R{落}{お}ちていきます。\R{落}{お}ちていくようすを\R{友}{とも}だちに\R{伝}{つた}えるとしたら、どう\R{言}{い}うかな？

\par\vspace{1mm}
\noindent\begin{minipage}[c]{0.36\linewidth}\centering
\begin{tikzpicture}[x=1mm, y=1mm]
  % 1秒ごとのボールの位置（落ちた距離 4.9t^2 m を 1.3 mm/m で）。速さの矢印は 0.55 mm（m/s）
  \draw[bkLine, line width=0.6pt] (16,3) -- (16,-76);
  \node[anchor=east, font=\sffamily\small, text=bkInk, inner sep=0pt] at (10,1.6) {0\R{秒}{びょう}};
  \node[anchor=east, font=\sffamily\small, text=bkInk, inner sep=0pt] at (10,-7.6) {1\R{秒後}{びょう|ご}};
  \node[anchor=east, font=\sffamily\small, text=bkInk, inner sep=0pt] at (10,-25.5) {2\R{秒後}{びょう|ご}};
  \node[anchor=east, font=\sffamily\small, text=bkOrange!80!black, inner sep=0pt] at (10,-57.3) {3\R{秒後}{びょう|ご}};
  \bkiiaball{(16,0)}{2.4}{0}
  \bkiiaball{(16,-6.4)}{2.4}{0}
  \bkiiaball{(16,-25.5)}{2.4}{0}
  \bkiiaball{(16,-57.3)}{2.4}{1}
  % 速さの矢印
  \draw[bk arrow, bkBlue] (24,-6.4) -- ++(0,-5.5);
  \draw[bk arrow, bkBlue] (24,-25.5) -- ++(0,-11);
  \draw[bk arrow, bkBlue, dash pattern=on 3pt off 2pt] (24,-57.3) -- ++(0,-16.5);
  \node[anchor=west, font=\sffamily\small, text=bkInk, inner sep=0pt] at (27,0) {\R{速}{はや}さ 0\,m/s};
  \node[anchor=west, font=\sffamily\small, text=bkBlue!80!black, inner sep=0pt] at (27,-9.6) {\R{約}{やく}\qty{10}{m/s}};
  \node[anchor=west, font=\sffamily\small, text=bkBlue!80!black, inner sep=0pt] at (27,-31) {\R{約}{やく}\qty{20}{m/s}};
  \node[anchor=west, font=\sffamily\bfseries\small, text=bkOrange!80!black, inner sep=0pt] at (27,-65.5) {？};
\end{tikzpicture}
\end{minipage}\hfill
\begin{minipage}[c]{0.61\linewidth}
\begin{tikzpicture}[x=1mm, y=1mm]
  \node[anchor=north west, draw=bkGray, line width=0.9pt, rounded corners=2mm, fill=bkPaper, text width=74mm,
    inner xsep=3.5mm, inner ysep=2.5mm, align=left, font=\small] (k) at (0,0)
    {{\sffamily\bfseries\color{bkGray}ことばで\R{言}{い}うと}\\
     「ボールは、だんだん\R{速}{はや}くなる。」\\
     →\,「だんだん」が、どれくらいかわからない。\R{人}{ひと}によって、\R{思}{おも}いうかべる\R{速}{はや}さがちがう。};
  \node[anchor=north west, draw=bkTheme, line width=1.2pt, rounded corners=2mm, fill=white, text width=74mm,
    inner xsep=3.5mm, inner ysep=2.5mm, align=left, font=\small] (s) at ([yshift=-4mm]k.south west)
    {{\sffamily\bfseries\color{bkDeep}\R{数}{かず}と\R{単位}{たん|い}で\R{言}{い}うと}\\
     「ボールの\R{速}{はや}さは、1\R{秒}{びょう}ごとに\R{約}{やく}\qty{10}{m/s}ずつふえる。」\\
     →\,だれが\R{読}{よ}んでも、\R{同}{おな}じ\R{速}{はや}さを\R{思}{おも}いうかべられる。};
\end{tikzpicture}
\end{minipage}

\par\vspace{2mm}
「だんだん\R{速}{はや}くなる」は、まちがいではありません。でも、ゆっくり\R{速}{はや}くなるのか、ぐんぐん\R{速}{はや}くなるのか、ことばだけではわかりません。

\R{数}{かず}と\R{単位}{たん|い}で\R{言}{い}うと、「1\R{秒}{びょう}ごとに\R{約}{やく}\qty{10}{m/s}ずつふえる」です。\R{左}{ひだり}の\R{図}{ず}のように、はなした\R{瞬間}{しゅん|かん}の\R{速}{はや}さは \qty{0}{m/s}。1\R{秒後}{びょう|ご}には\R{約}{やく}\qty{10}{m/s}、2\R{秒後}{びょう|ご}には\R{約}{やく}\qty{20}{m/s}。1\R{秒}{びょう}たつたびに、\R{速}{はや}さに\R{約}{やく}\qty{10}{m/s}がたされていきます。

\par\vspace{1mm}
\noindent{\sffamily\bfseries\color{bkDeep}\qty{10}{m/s}って、どんな\R{速}{はや}さ？}\par
m/s は「メートル\R{毎秒}{まい|びょう}」と\R{読}{よ}みます（くわしくは2-3）。\qty{10}{m/s}は、1\R{秒間}{びょう|かん}に\qty{10}{m}\R{進}{すす}む\R{速}{はや}さです。\qty{25}{m}プールなら、1\R{秒}{びょう}で\qty{10}{m}、2\R{秒}{びょう}で\qty{20}{m}\R{進}{すす}み、3\R{秒}{びょう}たたないうちに\R{向}{む}こうのはしに\R{着}{つ}きます。

\begin{bkfigure}
\begin{tikzpicture}[x=1mm, y=1mm]
  % 25 m プール（1 m を 4.4 mm で）
  \fill[bkBlue!14] (0,0) rectangle (110,9);
  \draw[bkBlue!70!black, line width=0.8pt] (0,0) rectangle (110,9);
  \foreach \m/\t in {0/{0\R{秒}{びょう}}, 10/{1\R{秒後}{びょう|ご}}, 20/{2\R{秒後}{びょう|ご}}} {
    \draw[bkInk, line width=0.7pt] (\m*4.4,-1.5) -- (\m*4.4,10.5);
    \node[anchor=south, font=\sffamily\footnotesize, text=bkInk, inner sep=0pt] at (\m*4.4,12) {\t};
    \node[anchor=north, font=\sffamily\footnotesize, text=bkInk, inner sep=0pt] at (\m*4.4,-3) {\m\,m};}
  \node[anchor=north, font=\sffamily\footnotesize, text=bkInk, inner sep=0pt] at (110,-3) {\qty{25}{m}};
  \foreach \m in {0,10} {\draw[bk arrow, bkOrange, line width=1.1pt] (\m*4.4+2,4.5) -- (\m*4.4+42,4.5);}
  \node[anchor=center, font=\sffamily\footnotesize, text=bkOrange!80!black, fill=bkBlue!14, inner sep=0.6mm] at (22,4.5) {\qty{10}{m}};
  \node[anchor=center, font=\sffamily\footnotesize, text=bkOrange!80!black, fill=bkBlue!14, inner sep=0.6mm] at (66,4.5) {\qty{10}{m}};
\end{tikzpicture}
\bkfigcaption{\qty{10}{m/s}で\R{進}{すす}むと、1\R{秒}{びょう}ごとに\qty{10}{m}ずつ\R{進}{すす}む。}
\end{bkfigure}

50\,m\R{走}{そう}を\R{約}{やく}8\R{秒}{びょう}で\R{走}{はし}る\R{人}{ひと}は、1\R{秒間}{びょう|かん}に\R{約}{やく}\qty{6}{m}\R{進}{すす}みます（$50\div8=6.25$）。\qty{10}{m/s}は、それよりずっと\R{速}{はや}い\R{速}{はや}さです。

\par\vspace{1mm}
\noindent{\sffamily\bfseries\color{bkDeep}なぜ「\R{約}{やく}10」なの？}\par
ボールの\R{落}{お}ち\R{方}{かた}は、\R{約}{やく}400\R{年前}{ねん|まえ}のイタリアの\R{科学者}{か|がく|しゃ}ガリレオのころから、たくさんの\R{人}{ひと}がくわしくはかってきました。はかると、\R{速}{はや}さは1\R{秒}{びょう}ごとに\R{約}{やく}\qty{9.8}{m/s}ずつふえます。この\R{本}{ほん}では、\R{計算}{けい|さん}しやすいように、9.8を\R{約}{やく}10として\R{使}{つか}います。「\R{約}{やく}」は「だいたい」という\R{意味}{い|み}です。

ただし、これは、\R{空気}{くう|き}がボールをおし\R{返}{かえ}す\R{力}{ちから}（\R{空気}{くう|き}の\R{抵抗}{てい|こう}）が\R{小}{ちい}さいときの\R{話}{はなし}です。\R{紙}{かみ}のように\R{軽}{かる}くて\R{広}{ひろ}いものは、\R{空気}{くう|き}におし\R{返}{かえ}されて、ゆっくり\R{落}{お}ちます。

%% ---------------- 数と単位で言うと、何がよい？ ----------------
\bksection{me}{\R{数}{かず}と\R{単位}{たん|い}で\R{言}{い}うと、\R{何}{なに}がよい？}

ことばを\R{数}{かず}と\R{単位}{たん|い}で\R{言}{い}いかえると、よいことが3つあります。

\bksubsection{1}{くわしく\R{伝}{つた}わる}

\begin{bkcenter}
\begin{bktable}{colspec={Q[l,m,wd=44mm] Q[l,m,wd=64mm]}, rowsep=1.5pt, cells={font=\small}, row{1}={font=\sffamily\bfseries\small}}
ことばで\R{言}{い}うと & \R{数}{かず}と\R{単位}{たん|い}で\R{言}{い}うと \\
\R{駅}{えき}まで、もうすぐ。 & \R{駅}{えき}まで、\R{歩}{ある}いて5\R{分}{ふん}。 \\
\R{今日}{きょう}は\R{暑}{あつ}い。 & \R{今日}{きょう}の\R{最高気温}{さい|こう|き|おん}は \qty{35}{\degreeCelsius}。 \\
\R{重}{おも}い\R{荷物}{に|もつ}。 & \R{重}{おも}さ \qty{10}{kg} の\R{荷物}{に|もつ}。 \\
ボールが、だんだん\R{速}{はや}くなる。 & \R{速}{はや}さが、1\R{秒}{びょう}ごとに\R{約}{やく}\qty{10}{m/s}ずつふえる。 \\
\end{bktable}
\bktabcaption{ことばと、\R{数}{かず}と\R{単位}{たん|い}}
\end{bkcenter}

「\R{暑}{あつ}い」と\R{感}{かん}じるかどうかは、\R{人}{ひと}によってちがいます。でも、「\qty{35}{\degreeCelsius}」は、だれが\R{温度計}{おん|ど|けい}を\R{読}{よ}んでも \qty{35}{\degreeCelsius} です。\R{数}{かず}と\R{単位}{たん|い}は、\R{人}{ひと}によって\R{変}{か}わらない「ものさし」ではかった\R{結果}{けっ|か}なので、くわしく、\R{正}{ただ}しく\R{伝}{つた}わります。

\bksubsection{2}{くらべられる}

\R{月}{つき}の\R{上}{うえ}でボールをはなしても、ボールは、だんだん\R{速}{はや}くなりながら\R{落}{お}ちます。ことばで\R{言}{い}うと、\R{地球}{ち|きゅう}とまったく\R{同}{おな}じです。でも、はかってみると、\R{月}{つき}では\R{速}{はや}さが1\R{秒}{びょう}ごとに\R{約}{やく}\qty{1.6}{m/s}ずつしかふえません。

\noindent\begin{minipage}[c]{0.52\linewidth}\centering
\begin{tikzpicture}[x=1mm, y=1mm]
  % 速さの棒グラフ（1 m/s を 1.1 mm で）
  \draw[bkInk, line width=0.7pt] (0,0) -- (62,0);
  \foreach \t/\e/\m/\ml in {1/10/1.6/1.6, 2/20/3.2/3.2, 3/30/4.8/4.8} {
    \fill[bkOrange] (\t*19-13,0) rectangle ++(7,\e*1.1);
    \fill[bkBlue!55] (\t*19-5,0) rectangle ++(7,\m*1.1);
    \node[anchor=south, font=\sffamily\footnotesize, text=bkOrange!80!black, inner sep=0pt] at (\t*19-9.5,\e*1.1+1.5) {\e};
    \node[anchor=south, font=\sffamily\footnotesize, text=bkBlue!80!black, inner sep=0pt] at (\t*19-1.5,\m*1.1+1.5) {\ml};
    \node[anchor=north, font=\sffamily\footnotesize, text=bkInk, inner sep=0pt] at (\t*19-5.5,-2) {\t\R{秒後}{びょう|ご}};}
  \fill[bkOrange] (0,46) rectangle ++(4,3);
  \node[anchor=west, font=\sffamily\footnotesize, inner sep=0pt] at (5.5,47.5) {\R{地球}{ち|きゅう}};
  \fill[bkBlue!55] (20,46) rectangle ++(4,3);
  \node[anchor=west, font=\sffamily\footnotesize, inner sep=0pt] at (25.5,47.5) {\R{月}{つき}};
  \node[anchor=west, font=\sffamily\footnotesize, text=bkGray, inner sep=0pt] at (36,47.5) {（\R{単位}{たん|い}は m/s）};
\end{tikzpicture}
\end{minipage}\hfill
\begin{minipage}[c]{0.45\linewidth}
\R{同}{おな}じ「だんだん\R{速}{はや}くなる」でも、\R{数}{かず}でくらべると、\R{地球}{ち|きゅう}のボールのほうが、ずっと\R{速}{はや}く\R{速}{はや}くなることがわかります。

1.6を6\R{倍}{ばい}すると9.6で、ほぼ10。\R{地球}{ち|きゅう}でふえる\R{速}{はや}さは、\R{月}{つき}の\R{約}{やく}6\R{倍}{ばい}です。\R{月}{つき}は\R{地球}{ち|きゅう}より\R{小}{ちい}さく、ものを\R{引}{ひ}っぱる\R{力}{ちから}が\R{弱}{よわ}いからです。
\end{minipage}

\bksubsection{3}{\R{計算}{けい|さん}できる}

\R{地球}{ち|きゅう}のボールの、3\R{秒後}{びょう|ご}の\R{速}{はや}さはどれくらいでしょう。\R{数}{かず}と\R{単位}{たん|い}で\R{書}{か}いてあれば、\R{計算}{けい|さん}で\R{求}{もと}められます。

\begin{bkfigure}
\begin{tikzpicture}[x=1mm, y=1mm]
  \foreach \i/\t/\v in {0/{0\R{秒}{びょう}}/{\qty{0}{m/s}}, 1/{1\R{秒後}{びょう|ご}}/{\R{約}{やく}\qty{10}{m/s}}, 2/{2\R{秒後}{びょう|ご}}/{\R{約}{やく}\qty{20}{m/s}}} {
    \node[bk node, minimum width=24mm, inner ysep=1.5mm] (b\i) at (\i*36,0) {\t\\{\bfseries \v}};}
  \node[bk node, draw=bkOrange, line width=1.4pt, minimum width=24mm, inner ysep=1.5mm] (b3) at (108,0)
    {3\R{秒後}{びょう|ご}\\{\bfseries\color{bkOrange!80!black}\R{約}{やく}\qty{30}{m/s}}};
  \foreach \a/\b in {b0/b1,b1/b2,b2/b3} \draw[bk link] (\a) -- node[above, bk link label] {$+10$} (\b);
\end{tikzpicture}
\bkfigcaption{1\R{秒}{びょう}たつごとに、\qty{10}{m/s}ずつたしていく。}
\end{bkfigure}

3\R{秒}{びょう}のあいだに、\qty{10}{m/s}が3\R{回}{かい}たされます。だから、$0+10+10+10=30$ です。10が3つ\R{分}{ぶん}なので、かけ\R{算}{ざん}で $10\times3=30$ と\R{計算}{けい|さん}することもできます（1-12：かけ\R{算}{ざん}は\R{同}{おな}じ\R{数}{かず}のまとまり）。3\R{秒後}{びょう|ご}の\R{速}{はや}さは、\R{約}{やく}\qty{30}{m/s}です。

ここでも「\R{約}{やく}」がつくのは、ほんとうにふえる\R{速}{はや}さが1\R{秒}{びょう}ごとに\R{約}{やく}\qty{9.8}{m/s}だからです。$9.8\times3=29.4$ なので、ほんとうは\R{約}{やく}\qty{29.4}{m/s}。だいたい30です。

\par\vspace{1mm}
\noindent\begin{minipage}[c]{0.70\linewidth}
\qty{30}{m/s}は、どれくらいの\R{速}{はや}さでしょう。1\R{時間}{じ|かん}は60\R{分}{ぷん}、1\R{分}{ぷん}は60\R{秒}{びょう}なので、1\R{時間}{じ|かん}は $60\times60=3600$\R{秒}{びょう}です。1\R{秒}{びょう}に\qty{30}{m}\R{進}{すす}むと、1\R{時間}{じ|かん}（3600\R{秒}{びょう}）では $30\times3600=108000$ で、\qty{108000}{m}。\qty{1}{km}は\qty{1000}{m}なので、\qty{108}{km}です。\R{時速}{じ|そく}\qty{108}{km}で、\R{高速道路}{こう|そく|どう|ろ}を\R{走}{はし}る\R{車}{くるま}くらいの\R{速}{はや}さです（\R{時速}{じ|そく}と\R{秒速}{びょう|そく}の\R{換算}{かん|さん}は、2-10でくわしく\R{学}{まな}びます）。
\end{minipage}\hfill
\begin{minipage}[c]{0.27\linewidth}\centering
\begin{tikzpicture}[x=1mm, y=1mm]
  \bkcar{(3,0)}{1.1}{bkBlue}
  \node[anchor=north, font=\sffamily\footnotesize, text=bkInk, inner sep=0pt] at (16,-3) {\R{時速}{じ|そく}\R{約}{やく}\qty{108}{km}};
\end{tikzpicture}
\end{minipage}

\par\vspace{1mm}
\R{月}{つき}のボールなら、$1.6\times3=4.8$ で、3\R{秒後}{びょう|ご}は\R{約}{やく}\qty{4.8}{m/s}。これも、\R{計算}{けい|さん}で\R{求}{もと}められます。

\R{地球}{ち|きゅう}のボールにもどると、4\R{秒後}{びょう|ご}は $10\times4=40$ で\R{約}{やく}\qty{40}{m/s}、5\R{秒後}{びょう|ご}は $10\times5=50$ で\R{約}{やく}\qty{50}{m/s}。「10×\R{秒数}{びょう|すう}」という\R{同}{おな}じ\R{計算}{けい|さん}で、\R{何秒後}{なん|びょう|ご}の\R{速}{はや}さでも\R{求}{もと}められます（\R{空気}{くう|き}の\R{抵抗}{てい|こう}が\R{小}{ちい}さいあいだ）。

\begin{taisetsu}
ことばを\R{数}{かず}と\R{単位}{たん|い}で\R{言}{い}いかえると、\\
①\,くわしく\R{伝}{つた}わる\quad ②\,くらべられる\quad ③\,\R{計算}{けい|さん}できる
\end{taisetsu}

\bkfukidashi{mina}{「すぐ」「たくさん」「とても\R{重}{おも}い」も、\R{数}{かず}と\R{単位}{たん|い}で\R{言}{い}いかえられるかな？\quad この\R{巻}{かん}で、そのことばを1つずつ\R{学}{まな}ぼう。}


\clearpage
\bksection{card}{この\R{巻}{かん}で\R{学}{まな}ぶ「\R{数}{かず}と\R{量}{りょう}のことば」}

ふだんのことばを\R{数}{かず}と\R{単位}{たん|い}で\R{言}{い}いかえるには、いろいろな\R{道具}{どう|ぐ}が\R{必要}{ひつ|よう}です。この\R{巻}{かん}の\R{章}{しょう}で、1つずつ\R{学}{まな}んでいきます。

\begin{bkcenter}
\begin{bktable}{colspec={Q[l,m,wd=50mm] Q[l,m,wd=56mm] Q[c,m,wd=14mm]}, rowsep=1.2pt, cells={font=\small}, row{1}={font=\sffamily\bfseries\small}}
ことばで\R{言}{い}うと & \R{数}{かず}と\R{単位}{たん|い}で\R{言}{い}うと & \R{学}{まな}ぶ\R{章}{しょう} \\
\R{駅}{えき}まで、あと5。 & \R{駅}{えき}まで、あと5\R{分}{ふん}（\R{単位}{たん|い}をつける） & 2-1 \\
\R{今朝}{け|さ}は、とても\R{寒}{さむ}い。 & \R{今朝}{け|さ}の\R{気温}{き|おん}は \qty{-3}{\degreeCelsius} & 2-2 \\
ピザを3\R{人}{にん}で\R{同}{おな}じに\R{分}{わ}ける。 & 1\R{人}{にん}\R{分}{ぶん}は $\dfrac{1}{3}$\R{枚}{まい} & 2-4 \\
いつもより\R{少}{すこ}し\R{安}{やす}い。 & 20\,\%\R{引}{び}き（0.8\R{倍}{ばい}の\R{値段}{ね|だん}） & 2-5 \\
だいたい5\,cm くらい。 & \qty{5.0}{cm}（ずれは\R{約}{やく}\qty{0.05}{cm}） & 2-6 \\
\R{光}{ひかり}は、ものすごく\R{速}{はや}い。 & \R{秒速}{びょう|そく}\R{約}{やく}$3.0\times10^{8}$\,m & 2-7 \\
どんな\R{速}{はや}さでも\R{使}{つか}えるきまり。 & \R{道}{みち}のり $x$＝\R{速}{はや}さ $v$×\R{時間}{じ|かん} $t$ & 2-11 \\
\end{bktable}
\bktabcaption{この\R{巻}{かん}で\R{学}{まな}ぶ、ことばの\R{言}{い}いかえ}
\end{bkcenter}

\begin{yattemiyou}[\R{身}{み}のまわりのことばを、\R{数}{かず}と\R{単位}{たん|い}で\R{言}{い}いかえよう]
\begin{enumerate}[label={\bkcircled[bkGreen]{\arabic*}}, leftmargin=2.4\zw, itemsep=1pt, topsep=2pt]
\item 「\R{重}{おも}いランドセル」：\R{体重計}{たい|じゅう|けい}で、ランドセルを\R{持}{も}ったときと\R{持}{も}たないときの\R{体重}{たい|じゅう}をはかり、ひき\R{算}{ざん}で\R{何}{なん}kg かを\R{求}{もと}めよう。
\item 「\R{近}{ちか}くの\R{公園}{こう|えん}」：\R{家}{いえ}から\R{歩}{ある}いて\R{何分}{なん|ぷん}かかるかを、\R{時計}{と|けい}ではかろう。
\item 「\R{今日}{きょう}は\R{暑}{あつ}い」：\R{天気予報}{てん|き|よ|ほう}の\R{最高気温}{さい|こう|き|おん}を\R{調}{しら}べて、\R{昨日}{きのう}とくらべよう。\R{何}{なん}℃ ちがうかな？
\end{enumerate}
ノートに「ことば → \R{数}{かず}と\R{単位}{たん|い}」の\R{表}{ひょう}を\R{作}{つく}って\R{書}{か}きこもう。\R{答}{こた}えには、かならず\R{単位}{たん|い}をつけます。
\end{yattemiyou}


\clearpage
%% ---------------- p.9 この巻が通る路線と駅 ----------------
\bksection{bkiiaroute}{この\R{巻}{かん}が\R{通}{とお}る\R{路線}{ろ|せん}と\R{駅}{えき}}

この\R{本}{ほん}の\R{章}{しょう}は、7\R{本}{ほん}の\R{路線}{ろ|せん}の\R{駅}{えき}です（4・5ページ）。\R{同}{おな}じ\R{路線}{ろ|せん}の\R{駅}{えき}は、1つの\R{考}{かんが}え\R{方}{かた}でつながっています。

\begin{bkfigure}
\begin{tikzpicture}[x=1mm, y=1mm]
  \bkiiaroute{70}{1}{1-5}{{\bfseries 2-3}・2-4・2-15・2-16・2-19}{3-5・5-1}
  \bkiiaroute{56}{2}{1-21}{{\bfseries 2-3}}{4-5・5-10}
  \bkiiaroute{42}{3}{ここが\R{始発}{し|はつ}}{{\bfseries 2-7}・2-8・2-17}{3-11・5-14}
  \bkiiaroute{28}{5}{1-13}{2-11〜2-14・{\bfseries 2-18}}{3-6・7-1}
  \bkiiaroute{14}{6}{ここが\R{始発}{し|はつ}}{{\bfseries 2-2}・2-9}{3-2・4-10}
  \bkiiaroute{0}{7}{1-11}{2-1・2-5・{\bfseries 2-6}・2-10}{3-14}
  \node[font=\sffamily\bfseries\small, text=bkGray] at (31,81) {\R{前}{まえ}の\R{巻}{かん}};
  \node[font=\sffamily\bfseries\small, text=bkDeep] at (75,81) {この\R{巻}{かん}（\R{第}{だい}2\R{巻}{かん}）};
  \node[font=\sffamily\bfseries\small, text=bkGray] at (119,81) {この\R{先}{さき}};
\end{tikzpicture}
\bkfigcaption{\R{第}{だい}2\R{巻}{かん}を\R{通}{とお}る6\R{本}{ぽん}の\R{路線}{ろ|せん}。\R{四角}{し|かく}の\R{中}{なか}がこの\R{巻}{かん}の\R{駅}{えき}（\R{太字}{ふと|じ}は\R{大}{おお}きな\R{駅}{えき}）。}
\end{bkfigure}

\begin{bkcenter}
\begin{bktable}{colspec={Q[c,m,wd=7mm] Q[l,m,wd=19mm] Q[l,m,wd=98mm]}, rowsep=1.5pt, cells={font=\small}, row{1}={font=\sffamily\bfseries\small}}
 & \R{路線}{ろ|せん} & \R{大}{おお}きな\R{駅}{えき}で\R{学}{まな}ぶこと \\
\bkrouteicon{1} & \R{変化}{へん|か} & ひき\R{算}{ざん}で\R{変}{か}わった\R{量}{りょう}を\R{求}{もと}める（2-3）→ \R{傾}{かたむ}き（3-5） \\
\bkrouteicon{2} & たし\R{合}{あ}わせ & かけ\R{算}{ざん}の\R{答}{こた}えを\R{面積}{めん|せき}で\R{見}{み}る（2-3）→ \R{積分}{せき|ぶん}（5-10） \\
\bkrouteicon{3} & \R{倍々}{ばい|ばい} & \R{同}{おな}じ\R{数}{かず}を\R{何回}{なん|かい}もかける（2-7）→ \R{倍々}{ばい|ばい}にふえる\R{量}{りょう}（3-11） \\
\bkrouteicon{5} & \R{表}{ひょう} & \R{表計算}{ひょう|けい|さん}で、くり\R{返}{かえ}し\R{計算}{けい|さん}する（2-18）→ \R{連立方程式}{れん|りつ|ほう|てい|しき}（3-6） \\
\bkrouteicon{6} & \R{矢印}{や|じるし} & マイナスは\R{反対}{はん|たい}の\R{向}{む}き（2-2）→ \R{座標}{ざ|ひょう}（3-2）→ ベクトル \\
\bkrouteicon{7} & はかる & はかった\R{数}{かず}の「ずれ」（2-6）→ ずれを\R{数}{かず}で\R{表}{あらわ}す（3-14） \\
\end{bktable}
\end{bkcenter}

\R{倍々}{ばい|ばい}の\R{路線}{ろ|せん}と\R{矢印}{や|じるし}の\R{路線}{ろ|せん}は、この\R{巻}{かん}が\R{始発駅}{し|はつ|えき}です。

\clearpage
%% ---------------- p.10 この巻でできるようになること ----------------
\bksection{goal}{この\R{巻}{かん}でできるようになること}

\R{第}{だい}2\R{巻}{かん}は、4つの\R{部}{ぶ}に\R{分}{わ}かれています。

\begin{bkfigure}
\begin{tikzpicture}[x=1mm, y=1mm,
  pt/.style={anchor=north west, draw=bkTheme, line width=1pt, fill=white, rounded corners=2.5mm,
    minimum width=65mm, minimum height=52mm, inner sep=0pt},
  pn/.style={anchor=north west, fill=bkTheme, text=bkOn, rounded corners=1.2mm, font=\sffamily\bfseries\small,
    inner xsep=2mm, inner ysep=1mm},
  pd/.style={anchor=north west, text width=58mm, align=left, font=\small, inner sep=0pt},
  px/.style={anchor=north, text=bkDeep, inner sep=0pt}]
  \foreach \x/\y in {0/0, 69/0, 0/-56, 69/-56} {\node[pt] at (\x,\y) {};}
  % 第1部
  \node[pn] at (3.5,-3.5) {\R{第}{だい}1\R{部}{ぶ}\enspace\R{量}{りょう}と\R{数}{かず}};
  \node[px] at (32.5,-14) {$\qty{5}{m}=\qty{500}{cm}$\qquad$-3+5=2$};
  \node[pd] at (3.5,-23) {\bkiiaitem{\R{量}{りょう}を、\R{数}{かず}と\R{単位}{たん|い}で\R{書}{か}ける。}\bkiiaitem{マイナス・\R{分数}{ぶん|すう}・\R{割合}{わり|あい}を\R{使}{つか}って\R{計算}{けい|さん}できる。}\bkiiaitem{はかった\R{数}{かず}のずれがわかる。}};
  % 第2部
  \node[pn] at (72.5,-3.5) {\R{第}{だい}2\R{部}{ぶ}\enspace\R{大}{おお}きな\R{数}{かず}・\R{小}{ちい}さな\R{数}{かず}・\R{単位}{たん|い}};
  \node[px] at (101.5,-14) {$3.0\times10^{8}$\,m/s\qquad$\sqrt{2}=1.41\cdots$};
  \node[pd] at (72.5,-23) {\bkiiaitem{\R{大}{おお}きな\R{数}{かず}・\R{小}{ちい}さな\R{数}{かず}を、10の\R{何乗}{なん|じょう}で\R{書}{か}ける。}\bkiiaitem{2\R{回}{かい}かけて2になる\R{数}{かず}がわかる。}\bkiiaitem{\R{単位}{たん|い}をそろえて\R{計算}{けい|さん}できる。}};
  % 第3部
  \node[pn] at (3.5,-59.5) {\R{第}{だい}3\R{部}{ぶ}\enspace\R{文字}{も|じ}と\R{式}{しき}};
  \node[px] at (32.5,-68.5) {$x=v\times t$\enspace→\enspace$t=x\div v$};
  \node[pd] at (3.5,-79) {\bkiiaitem{\R{文字}{も|じ}を\R{使}{つか}って、きまりを1\R{行}{ぎょう}で\R{書}{か}ける。}\bkiiaitem{\R{式}{しき}を\R{変形}{へん|けい}して、ほしい\R{量}{りょう}をとり\R{出}{だ}せる。}};
  % 第4部
  \node[pn] at (72.5,-59.5) {\R{第}{だい}4\R{部}{ぶ}\enspace コンピュータで\R{計算}{けい|さん}};
  \node[px] at (101.5,-68.5) {$1011_{(2)}=11$};
  \node[pd] at (72.5,-79) {\bkiiaitem{0と1だけで\R{数}{かず}を\R{書}{か}ける。}\bkiiaitem{\R{表計算}{ひょう|けい|さん}とプログラムで、くり\R{返}{かえ}しの\R{計算}{けい|さん}ができる。}};
\end{tikzpicture}
\bkfigcaption{4つの\R{部}{ぶ}と、できるようになること}
\end{bkfigure}

\begin{column}{なぜ「\R{数}{かず}と\R{量}{りょう}のことば」なの？}
「\qty{5}{m}」と\R{書}{か}けば、\R{日本}{に|ほん}でもブラジルでも、だれが\R{読}{よ}んでも\R{同}{おな}じ\R{長}{なが}さです。\R{数}{かず}・\R{単位}{たん|い}・\R{記号}{き|ごう}・\R{式}{しき}は、\R{世界}{せ|かい}じゅうの\R{人}{ひと}が、\R{物}{もの}の\R{動}{うご}きや\R{機械}{き|かい}のしくみを\R{正確}{せい|かく}に\R{伝}{つた}え\R{合}{あ}うための、\R{共通}{きょう|つう}のことばです。\R{第}{だい}3\R{巻}{かん}では、このことばを\R{使}{つか}って、\R{表}{ひょう}・グラフ・\R{式}{しき}から\R{変化}{へん|か}のきまりを\R{見}{み}つけます。
\end{column}

\clearpage
%% ---------------- p.11 この巻の読み方 ----------------
\bksection{hon}{この\R{巻}{かん}の\R{読}{よ}み\R{方}{かた}}

\begin{bkfigure}
\begin{tikzpicture}[x=1mm, y=1mm]
  \bkiiatalk{0}{sora}{ソラ}{\R{第}{だい}1\R{巻}{かん}では、\R{数}{かぞ}えたり、はかったりしたね。\R{第}{だい}2\R{巻}{かん}では\R{何}{なに}をするの？}
  \bkiiatalk{-19}{mina}{ミナ}{0より\R{小}{ちい}さい\R{数}{かず}や、\R{文字}{も|じ}の $x$ も\R{出}{で}てくるんだって。むずかしそう……。}
  \bkiiatalk{-38}{delta}{デルタ}{だいじょうぶ。どの\R{章}{しょう}も、\R{場面}{ば|めん}を\R{見}{み}て、\R{手}{て}を\R{動}{うご}かして、\R{気}{き}づいてから\R{式}{しき}にするよ。わからなくなったら、\R{前}{まえ}の\R{章}{しょう}にもどればいいよ。}
\end{tikzpicture}
\end{bkfigure}

\begin{bkfigure}
\begin{tikzpicture}[x=1mm, y=1mm]
  \bkiiacard{(0,0)}{omoidasu}{bkGray}{①\enspace\R{思}{おも}い\R{出}{だ}そう}
    {\R{章}{しょう}のはじめに、\R{前}{まえ}に\R{学}{まな}んだことを1\R{行}{ぎょう}で\R{思}{おも}い\R{出}{だ}します。\\[2pt]
     \bkiiapill{1-8}\enspace\qty{1}{m}＝\qty{100}{cm}\\[2pt]
     \R{札}{ふだ}の\R{章}{しょう}にもどれば、くわしく\R{書}{か}いてあります。}
  \bkiiacard{(70,0)}{doguu}{bkNavy}{②\enspace\R{電卓}{でん|たく}を\R{使}{つか}おう}
    {2-1から\R{使}{つか}ってかまいません。\\
     \bkiiaitem{おす\R{前}{まえ}に、\R{式}{しき}をノートに\R{書}{か}く。}\bkiiaitem{\R{答}{こた}えのだいたいの\R{大}{おお}きさが\R{合}{あ}っているか、\R{確}{たし}かめる。}}
  \bkiiacard{(0,-51)}{jibun}{bkOrange}{③\enspace ノートの\R{使}{つか}い\R{方}{かた}}
    {\bkiiaitem{\R{式}{しき}は1\R{行}{ぎょう}に1つ。「＝」をたてにそろえる。}\bkiiaitem{\R{答}{こた}えには\R{単位}{たん|い}をつける。}\bkiiaitem{\R{式}{しき}の\R{前}{まえ}に、\R{図}{ず}や\R{表}{ひょう}をかく。}}
  \bkiiacard{(70,-51)}{reidai}{bkBlue}{④\enspace\R{例題}{れい|だい}から\R{章末問題}{しょう|まつ|もん|だい}へ}
    {\R{例題}{れい|だい}を\R{読}{よ}む → じぶんでとこう（\R{同}{おな}じ\R{形}{かたち}）→ \R{章末問題}{しょう|まつ|もん|だい}。\\[2pt]
     わからなかったら、\R{同}{おな}じ\R{形}{かたち}の\R{例題}{れい|だい}にもどろう。}
\end{tikzpicture}
\end{bkfigure}

まちがえても、だいじょうぶ。ノートの\R{式}{しき}を1\R{行}{ぎょう}ずつ\R{見}{み}なおして、どの\R{行}{ぎょう}でまちがえたかを\R{見}{み}つければいいのです。
